% ─────────────────────────────────────────────────────────────────────────────
% Deliverable 2 — Software Design Documentation
% MATH5320 Portfolio Risk Management System
% Columbia University, MATH GR 5320 Financial Risk Management, Spring 2026
% ─────────────────────────────────────────────────────────────────────────────
\documentclass[11pt]{article}
\usepackage[margin=1in]{geometry}
\usepackage[T1]{fontenc}
\usepackage[utf8]{inputenc}
\usepackage{lmodern}
\usepackage{microtype}
\usepackage{xcolor}
\usepackage{amsmath,amssymb}
\usepackage{booktabs,longtable,array,ragged2e}
\usepackage{graphicx}
\usepackage{float}
\usepackage{hyperref}
\usepackage{enumitem}
\usepackage{parskip}
\usepackage{fancyvrb}

\hypersetup{
  colorlinks=true,
  linkcolor=blue,
  urlcolor=blue,
  citecolor=blue,
  pdftitle={Software Design Documentation -- MATH5320},
  pdfauthor={Nigel Li, Michael Adegbite, Stella}
}

\setlist[itemize]{topsep=3pt,itemsep=2pt,parsep=0pt,leftmargin=1.6em}
\setlist[enumerate]{topsep=3pt,itemsep=2pt,parsep=0pt,leftmargin=1.6em}
\newcolumntype{L}[1]{>{\RaggedRight\arraybackslash}p{#1}}
\newcommand{\code}[1]{\texttt{#1}}

% Compact verbatim for shell commands
\DefineVerbatimEnvironment{shellcode}{Verbatim}{fontsize=\small,xleftmargin=1em}

% ─────────────────────────────────────────────────────────────────────────────
\begin{document}

% ── Cover page ────────────────────────────────────────────────────────────────
\begin{titlepage}
\centering
\vspace*{2cm}
{\LARGE\bfseries Software Design Documentation\par}
\vspace{0.4cm}
{\large MATH GR 5320 Portfolio Risk Management System\par}
\vspace{0.4cm}
{\large Columbia University, Financial Risk Management, Spring 2026\par}
\vspace{1.4cm}

\begin{tabular}{L{4cm}L{9cm}}
\toprule
\textbf{Field} & \textbf{Value} \\
\midrule
Deliverable & 2 of 5 (15 points) \\
Authors & Nigel Li, Michael Adegbite, Stella \\
Reference Commit & \texttt{5841589} (main branch, May 2026) \\
Submission Date & May 12, 2026 \\
Python Version & 3.12.2 \\
Test Suite & 624 tests, 0 failures \\
Statement Coverage & 95\% \\
\bottomrule
\end{tabular}

\vspace{1.4cm}
\begin{minipage}{0.78\linewidth}
\small\itshape
This document describes the architecture, module design, data flow,
validation strategy, and known limitations of the MATH5320 risk engine.
It is structured to satisfy the pre-deployment model-risk management
documentation requirements discussed in Lecture 5.
\end{minipage}
\end{titlepage}

\tableofcontents
\clearpage

% ─────────────────────────────────────────────────────────────────────────────
\section{Executive Summary}

The \code{MATH5320} repository implements a modular Python and Streamlit
risk engine for portfolios of stocks and European options.  The design
separates user-interface code, data loading, portfolio representation,
pricing logic, risk models, orchestration, and validation tests into
distinct layers.  This separation is deliberate: it keeps quantitative
functions testable independently of the Streamlit interface, and treats
the UI as a thin presentation layer over a reusable analytical core.

The system workflow (portfolio input, market data loading, risk
parameter configuration, risk analysis, and VaR backtesting) is
reflected directly in the codebase structure.  \code{app.py} controls
the top-level Streamlit flow; \code{src/ui/} contains panel-level logic;
\code{src/services/} orchestrates end-to-end computations;
\code{src/portfolio/} and \code{src/schemas.py} define the domain model;
\code{src/pricing/} and \code{src/risk/} hold the analytical core; and
\code{tests/} provides broad validation coverage.

The design satisfies all five Lecture~5 pre-deployment requirements:
clear statement of purpose, design documentation, data analysis,
testing, and system analysis and testing.  The software architecture is
well suited to a model-risk-sensitive context precisely because it keeps
analytical logic inspectable, testable, and independent of the
user-interface layer.

% ─────────────────────────────────────────────────────────────────────────────
\section{System Purpose and Scope}

\subsection{Purpose}

The system is an educational portfolio risk application for MATH GR~5320.
Its core purpose is to enable a user to:

\begin{itemize}
  \item define a portfolio of stocks and European options;
  \item load historical market data via CSV upload or Yahoo Finance;
  \item compute Value at Risk and Expected Shortfall under historical
    simulation, delta-normal parametric, and Monte Carlo methods;
  \item choose historical or manual calibration of market-risk parameters;
  \item compare methods under common data and parameter assumptions;
  \item run walk-forward VaR backtests with Kupiec and Christoffersen diagnostics;
  \item inspect supporting analytics and download structured outputs.
\end{itemize}

The repository also contains extension modules for exact GBM/lognormal
risk, hazard-rate credit, Merton structural default, CDS, CVA,
counterparty mitigation, and illustrative regulatory capital calculations.
These are explicitly documented as course-formula extensions, not
production supervisory models.

\subsection{Scope}

\textbf{In scope:}
stocks; European calls and puts; historical, parametric, and Monte Carlo
VaR and ES; walk-forward VaR backtesting with Kupiec and Christoffersen
diagnostics; CSV and Yahoo Finance data loading; downloadable risk outputs;
and course-formula validation modules.

\textbf{Out of scope:}
production trading or risk management; enterprise authorization, audit,
and access control; American or path-dependent option pricing; full
volatility-surface or stochastic-volatility repricing; official
supervisory DFAST/CCAR modeling; and production XVA systems.

\subsection{Lecture~5 Design Compliance Matrix}

\begin{center}
\begin{tabular}{L{6.5cm}L{7.5cm}}
\toprule
\textbf{Lecture~5 / model-risk design requirement} & \textbf{Where this report satisfies it} \\
\midrule
Clear statement of purpose        & Section~2, System Purpose and Scope \\
Design documentation              & Sections 3–5 with diagrams and module inventories \\
Data analysis                     & Sections 4, 6, and 8 \\
Testing                           & Section 10 and cross-references to test files \\
System analysis and testing       & Sections 4, 7, and 10 \\
Module documentation              & Section~5 with purpose, inputs, outputs, test evidence \\
Data-flow documentation           & Section~4 with portfolio-to-output diagrams \\
Code review concerns              & Sections 8, 9, and 12 \\
Separation of concerns            & Sections 3 and 7 \\
\bottomrule
\end{tabular}
\end{center}

% ─────────────────────────────────────────────────────────────────────────────
\subsection{Model Choice Justification and Literature Basis}

Each method implemented in the repository is grounded in published research
or documented industry practice.

\begin{center}
\footnotesize
\begin{tabular}{L{2.8cm}L{4.8cm}L{5.5cm}}
\toprule
\textbf{Method} & \textbf{Justification} & \textbf{Reference} \\
\midrule
Black-Scholes pricing
  & Industry-standard European option pricing; BS delta feeds delta-normal VaR
  & Hull, \textit{Options, Futures, and Other Derivatives}, 11th ed., Chs.\ 15, 19 \\
Historical simulation
  & Non-parametric; no distributional assumption; captures fat tails from observed data
  & McNeil, Frey \& Embrechts, \textit{Quantitative Risk Management}, Princeton, 2015, Ch.~2 \\
Delta-normal parametric
  & Closed-form solution for linear portfolios; tractable baseline VaR
  & \textit{RiskMetrics Technical Document}, Morgan Guaranty Trust, 1994 \\
Monte Carlo VaR/ES
  & Full repricing for nonlinear payoffs; gold standard for options
  & McNeil et al.\ Ch.~2; Hull Ch.~22 \\
EWMA volatility
  & Down-weights old data; adapts faster to changing regimes than rolling equal-weight
  & \textit{RiskMetrics Technical Document}, 1994; Zivot \& Wang, Ch.~6 \\
Kupiec backtest
  & Likelihood-ratio test of unconditional exception frequency; the standard regulatory backtest
  & Kupiec (1995), \textit{FEDS Discussion Paper 95-24} \\
Christoffersen test
  & Conditional coverage test; detects exception clustering; required by Lecture~5
  & Christoffersen (1998), \textit{International Economic Review} 39(4) \\
\bottomrule
\end{tabular}
\end{center}

\subsection{Comparison with Alternative Approaches}

Three alternative implementations were considered and not adopted.  The
design choices are explicit, not accidental.

\begin{center}
\footnotesize
\begin{tabular}{L{2.8cm}L{3.0cm}L{3.2cm}L{4.0cm}}
\toprule
\textbf{Property} & \textbf{Historical simulation} & \textbf{Parametric} & \textbf{Monte Carlo} \\
\midrule
Distributional assumption & None (empirical) & Multivariate normal & Multivariate normal (extendable) \\
Option handling & Full repricing & Delta approximation & Full repricing \\
Computation & Moderate & Very fast, closed-form & Slow; scales with path count \\
Fat-tail capture & Yes, from data & No & Partial; depends on assumed distribution \\
Volatility clustering & None; equally weighted & Yes with EWMA & Yes with EWMA \\
Main limitation & Slow to adapt to regime changes & Underestimates nonlinear risk & Model-dependent; seed-sensitive \\
\bottomrule
\end{tabular}
\end{center}

Alternatives not implemented, with documented reasons:

\begin{itemize}
  \item \textbf{Age-weighted historical simulation.}  Assigns exponentially
    decaying weights to historical scenarios, addressing the slow-adaptation
    limitation.  Not implemented because the EWMA estimator on the
    parametric and MC engine achieves a comparable effect on covariance.
  \item \textbf{Extreme Value Theory (EVT).}  Parameterizes the tail of the
    loss distribution for superior tail estimation; documented as a future
    enhancement.
  \item \textbf{Filtered Historical Simulation (FHS).}  Applies GARCH
    filtering before historical simulation, directly addressing the
    Christoffersen exception-clustering failure observed in backtesting.
    Documented as the appropriate next step for a production build.
\end{itemize}

\subsection{Subjective Design Choices and Calibration Parameters}

Several parameter defaults required judgment.  These are documented here
rather than embedded silently in code defaults.

\begin{center}
\begin{tabular}{L{3.5cm}L{1.8cm}L{7.5cm}}
\toprule
\textbf{Parameter} & \textbf{Default} & \textbf{Justification} \\
\midrule
Lookback window
  & 252 days
  & One calendar year; standard RiskMetrics and Basel convention.  Shorter windows
    sacrifice stability; longer windows mix fundamentally different regimes. \\
EWMA half-life $N$
  & 60
  & Corresponds to $\lambda \approx 0.989$ decay factor per day;
    roughly quarterly effective memory.  Based on RiskMetrics recommendation
    for daily equity data. \\
Monte Carlo paths
  & 10\,000
  & Sufficient for 99\% VaR stability at academic scale.  A seeded test at
    100\,000 paths confirms convergence to within 2\% of exact GBM VaR;
    gains beyond 10\,000 are marginal for coursework purposes. \\
Vol floor
  & 1\%
  & Prevents numerical instability when historical vol approaches zero.
    Does not affect typical market conditions; applied only to option
    repricing vol-shock paths. \\
Confidence level (default)
  & 99\% VaR; 97.5\% ES
  & Aligns with Basel~III IMA: 99\% 10-day VaR for market risk reporting;
    97.5\% ES for FRTB purposes. \\
Risk horizon (default)
  & 5 days
  & Basel~III standard horizon for daily VaR reporting and backtesting. \\
\bottomrule
\end{tabular}
\end{center}

% ─────────────────────────────────────────────────────────────────────────────
\section{Software Architecture}

\subsection{Layered Architecture}

The system follows a strict layered architecture in which each layer has
a single responsibility and well-defined interfaces with adjacent layers:

\begin{center}
\fbox{\begin{minipage}{0.88\linewidth}
\centering
\textbf{User}\\[0.4em]
$\downarrow$\\[0.4em]
\textbf{UI Layer} \quad \code{app.py}, \code{src/ui/*}\\[0.4em]
$\downarrow$\\[0.4em]
\textbf{Service Layer} \quad \code{src/services/risk\_engine\_service.py}\\[0.4em]
$\downarrow$\\[0.4em]
\textbf{Domain Layer} \quad \code{src/schemas.py}, \code{src/portfolio/*}\\[0.4em]
$\downarrow$\\[0.4em]
\textbf{Model Layer} \quad \code{src/pricing/*}, \code{src/risk/*}, \code{src/credit/*}\\[0.4em]
$\downarrow$\\[0.4em]
\textbf{Outputs} \quad VaR/ES tables, loss distributions, backtest results, CSV/JSON downloads
\end{minipage}}
\end{center}

\subsection{Why This Architecture Fits a Risk Engine}

A layered architecture is appropriate for a model-risk-sensitive
codebase for several interconnected reasons.  Pure model functions can
be called directly from the test suite and from the Jupyter notebook
environment without any Streamlit state, which makes independent
validation straightforward.  Module responsibilities are narrow enough
to be read and understood in isolation.  Extension modules for credit,
CVA, and regulatory calculations can be added without modifying the
application shell or the core stock-and-option risk path.

\subsection{Separation of Concerns}

A key design decision is that the repository does not mix data
gathering, calibration, computation, and presentation within the same
module.  That discipline matters because model logic buried inside UI
callbacks or data-loading functions becomes difficult to validate and
govern.  The table below makes the separation explicit.

\begin{center}
\begin{tabular}{L{2.5cm}L{5.5cm}L{5.5cm}}
\toprule
\textbf{Layer} & \textbf{Responsibility} & \textbf{Must not do} \\
\midrule
Data layer      & Load, parse, clean, align, and validate prices & Compute VaR or ES \\
Estimator layer & Estimate means, volatilities, and covariances & Download data or render UI \\
Pricing layer   & Price stocks/options; compute option sensitivities & Portfolio-level workflow decisions \\
Portfolio layer & Aggregate positions, values, and exposures & Estimate risk-model parameters \\
Risk layer      & Compute VaR, ES, scenarios, and backtests & Render charts or collect user input \\
Service layer   & Orchestrate the end-to-end workflow & Hide quantitative formulas in UI code \\
UI layer        & Collect inputs and display outputs & Mutate model logic or silently change assumptions \\
\bottomrule
\end{tabular}
\end{center}

% ─────────────────────────────────────────────────────────────────────────────
\section{Data Flow and Control Flow}

\subsection{End-to-End Data-Flow Diagram}

\begin{center}
\fbox{\begin{minipage}{0.88\linewidth}
\centering
\textbf{Portfolio input} (stocks and European options)\\[0.35em]
$\downarrow$\\[0.35em]
\textbf{Input validation} (schemas and data checks)\\[0.35em]
$\downarrow$\\[0.35em]
\textbf{Market data loading} (CSV/Yahoo Finance, date alignment, missing-data checks)\\[0.35em]
$\downarrow$\\[0.35em]
\textbf{Current valuation} (stock value, option repricing, total portfolio value)\\[0.35em]
$\downarrow$\\[0.35em]
\textbf{Return and parameter estimation} (log returns, horizon returns, rolling/EWMA mean-covariance)\\[0.35em]
$\downarrow$\\[0.35em]
\textbf{Risk models:} \quad Historical VaR/ES \qquad Parametric VaR/ES \qquad Monte Carlo VaR/ES\\[0.35em]
$\downarrow$\\[0.35em]
\textbf{Outputs:} comparison table, losses CSV, JSON summary, plots, backtesting diagnostics
\end{minipage}}
\end{center}

\subsection{Core Modeling Conventions}

The software design follows the conventions documented in the README and
implemented in \code{src/risk/}:

\begin{itemize}
  \item Daily log returns: $r_{i,t} = \log(S_{i,t}/S_{i,t-1})$
  \item Overlapping horizon returns: $R_t^{(h)} = \sum_{k=0}^{h-1} r_{t-k}$
  \item Price shock: $S_{\mathrm{shocked}} = S_0 \cdot e^R$
  \item Portfolio PnL: $V_T - V_0$; loss: $V_0 - V_T$
  \item Horizon scaling: $\mu_h = \mu \cdot h$, $\Sigma_h = \Sigma \cdot h$
  \item Option pricing: Black-Scholes with continuous dividends
  \item Kupiec backtest statistic: $\mathrm{LR}_{\mathrm{uc}} \sim \chi^2_1$
\end{itemize}

These conventions define the data structures that flow between modules.
Historical and Monte Carlo engines consume shocked spot vectors and
full portfolio repricing, while the parametric engine consumes
delta-dollar exposure vectors and estimated covariance matrices.

\subsection{Backtesting Control-Flow Diagram}

\begin{center}
\fbox{\begin{minipage}{0.88\linewidth}
\centering
\textbf{Historical market data}\\[0.35em]
$\downarrow$\\[0.35em]
\textbf{Rolling estimation window:} use only data available at each forecast date\\[0.35em]
$\downarrow$\\[0.35em]
\textbf{VaR forecast}\\[0.35em]
$\downarrow$\\[0.35em]
\textbf{Realized future loss}\\[0.35em]
$\downarrow$\\[0.35em]
\textbf{Exception sequence:} record breach when realized loss exceeds VaR\\[0.35em]
$\downarrow$\\[0.35em]
\textbf{Diagnostics:} exception count, exception rate, Kupiec LR, Christoffersen LR, traffic-light zone
\end{minipage}}
\end{center}

Backtesting is implemented as an out-of-sample walk-forward process in
\code{src/risk/backtest.py}.  The implementation never uses future data
to estimate the VaR forecast at the current date: it slices prices up to
the forecast date, fits on that historical subset, and then compares the
resulting VaR with the realized future loss over the chosen horizon.

% ─────────────────────────────────────────────────────────────────────────────
\section{Module-by-Module Design}

\subsection{Module Inventory}

\begin{center}
\footnotesize
\begin{tabular}{L{2.4cm}L{3.8cm}L{2.4cm}L{2.2cm}L{2.4cm}}
\toprule
\textbf{Module} & \textbf{Purpose} & \textbf{Inputs} & \textbf{Outputs} & \textbf{Test evidence} \\
\midrule
Schemas & Define stock, option, and portfolio objects & User inputs & Structured portfolio objects & \code{test\_config\_and\_validation.py} \\
Market data & CSV and Yahoo Finance loading & Tickers, dates, CSVs & Aligned price data & \code{test\_market\_data.py} \\
Data validation & Validate price and input data & Prices, settings & Errors or clean acceptance & \code{test\_config\_and\_validation.py} \\
Black-Scholes & European option pricing and delta & $S,K,T,r,q,\sigma$, type & Price, delta & \code{test\_backend.py}, \code{test\_homework\_cases.py} \\
Position valuation & Per-position value, vol shock, delta-dollar sensitivity & Position and market inputs & Value, exposure & \code{test\_backend.py} \\
Portfolio valuation & Aggregate positions and exposures & Portfolio, spot vector & Total value, exposure vector & \code{test\_backend.py} \\
Returns & Log and horizon return construction & Price matrix & Return matrix & \code{test\_backend.py}, \code{test\_coverage\_gaps.py} \\
Estimators & Rolling, EWMA, and manual covariance assembly & Return matrix or manual params & Mean, covariance & \code{test\_backend.py}, \code{test\_homework\_cases.py} \\
Historical risk & Historical VaR/ES with full repricing and vol shock & Portfolio, history, settings & VaR, ES, losses & \code{test\_backend.py}, \code{test\_course\_validation.py} \\
Parametric risk & Delta-normal VaR/ES & Exposure vector, covariance & VaR, ES & \code{test\_backend.py}, \code{test\_es\_confidence\_split.py} \\
Monte Carlo risk & Simulated VaR/ES & Mean/cov, portfolio, vol-shock settings & VaR, ES, losses & \code{test\_backend.py}, \code{test\_coverage\_gaps.py} \\
Backtesting & Walk-forward VaR validation & History, model settings & Exceptions, Kupiec, Christoffersen & \code{test\_backend.py}, \code{test\_backtest\_extensions.py} \\
Exact GBM & Formula-sheet GBM VaR/ES & GBM parameters & Exact VaR, ES & \code{test\_lognormal.py}, \code{test\_course\_validation.py} \\
Hazard & Reduced-form default & Hazard, recovery, maturity & Survival, density, spread & \code{test\_credit.py}, \code{test\_course\_validation.py} \\
Merton & Structural default & $V_0,B,T,r,\mu,\sigma$ & PD, equity, debt & \code{test\_credit.py}, \code{test\_course\_validation.py} \\
CDS & Par spread calculation & Hazard, recovery, discounting & Spread, protection/premium values & \code{test\_credit.py}, \code{test\_course\_validation.py} \\
CVA & Counterparty valuation adjustment & Exposure, PD, recovery & CVA & \code{test\_credit.py}, \code{test\_cva\_mitigants.py} \\
Regulatory & RWA, capital ratio, DFAST-style calculations & Assets, losses, weights & Ratios, stress metrics & \code{test\_regulatory.py}, \code{test\_dfast\_pathing.py} \\
Service & Orchestrate end-to-end core risk run & Portfolio, data, settings & Unified result object & \code{test\_backend.py}, \code{integration\_test.py} \\
UI & Streamlit display and input logic & User interaction & Rendered panels & \code{test\_ui\_panels.py}, \code{test\_charts.py} \\
\bottomrule
\end{tabular}
\end{center}

\subsection{Public Interfaces and Contracts}

The main public interfaces define the contracts between the UI/service
layer and the model layer.

\begin{center}
\begin{tabular}{L{3.5cm}L{4.5cm}L{3.0cm}L{2.5cm}}
\toprule
\textbf{Interface} & \textbf{Inputs} & \textbf{Outputs} & \textbf{Failure mode} \\
\midrule
\code{black\_scholes\_price} & $S,K,T,\sigma,r,q$, option type & Option price & Reject invalid type or numerical domain \\
\code{black\_scholes\_delta} & Same inputs & Delta & Reject invalid type or numerical domain \\
\code{portfolio\_value} & Portfolio object, spot vector & Portfolio value & Fail visibly if required underlying missing \\
\code{portfolio\_delta\_dollar} & Portfolio, spot vector & Exposure vector & Fail visibly if option input inconsistent \\
\code{historical\_var\_es} & Portfolio, price history, horizon, confidences & VaR, ES, losses & Return explicit reason if history insufficient \\
\code{parametric\_var\_es} & Exposure vector, mean/covariance, horizon & VaR, ES & Reject invalid covariance or confidence \\
\code{monte\_carlo\_var\_es} & Portfolio, mean/cov, paths, seed & VaR, ES, losses & Reject invalid covariance or path count \\
\code{run\_backtest} & Portfolio, price history, model choice, lookback & Exception series, summary & Return explicit no-result reason when infeasible \\
\code{RiskEngineService} & Portfolio, prices, pricing date, settings & Unified result dict & Surface errors to the UI rather than swallowing \\
\bottomrule
\end{tabular}
\end{center}

\subsection{Core Data Structures}

\begin{center}
\begin{tabular}{L{3.5cm}L{5.5cm}L{4.5cm}}
\toprule
\textbf{Object} & \textbf{Main fields} & \textbf{Design role} \\
\midrule
\code{StockPosition}      & Ticker, quantity (signed) & Represents signed equity holdings \\
\code{OptionPosition}     & Underlying, type, strike, maturity, volatility, rate, dividend yield, multiplier, quantity & Represents European option contracts \\
\code{Portfolio}          & Stock positions, option positions & Single object passed through valuation and risk workflows \\
Risk settings bundle      & Lookback, horizon, confidences, estimator type, calibration mode, MC paths, vol-shock controls & Standardizes user configuration across risk engines \\
\code{ManualMarketParams} & Daily mean vector, volatility/covariance inputs & Allows parameter-driven parametric and MC runs \\
Risk result dictionary    & Method-level VaR, ES, losses, portfolio value, chart payloads & Feeds the UI and test artifacts \\
\code{BacktestResult}     & Forecast dates, realized losses, VaR forecasts, exceptions, diagnostics & Records outcome-analysis evidence \\
\bottomrule
\end{tabular}
\end{center}

\subsection{Operational Paths}

\textbf{Current risk run:}
\begin{enumerate}
  \item Validate portfolio and market-data availability.
  \item Compute current stock and option values.
  \item Estimate or accept market-risk parameters.
  \item Run historical, parametric, and Monte Carlo models.
  \item Aggregate VaR, ES, losses, and charts into one result object.
\end{enumerate}

\textbf{Backtest run:}
\begin{enumerate}
  \item Move through time using a rolling estimation window.
  \item Re-estimate the selected model at each date.
  \item Forecast VaR using only information available at that date.
  \item Observe the realized future loss.
  \item Record exceptions and compute Kupiec and Christoffersen diagnostics.
\end{enumerate}

\textbf{Formula-sheet extension run:}
\begin{enumerate}
  \item Call pure modules directly from notebooks or service helpers.
  \item Compare actual output with a deterministic course fixture.
  \item Assert equality within the documented numerical tolerance.
  \item Record the result in the validation artifacts.
\end{enumerate}

% ─────────────────────────────────────────────────────────────────────────────
\section{Input and Output Schemas}

Input enforcement is distributed across Streamlit UI controls,
\code{src/data/validation.py}, and numerical domain checks inside
pricing and model functions.  This section documents the intended
contract; specific enforcement layers are noted where relevant.

\subsection{Stock Input}

\begin{center}
\begin{tabular}{L{2.5cm}L{1.8cm}L{4cm}L{5cm}}
\toprule
\textbf{Field} & \textbf{Type} & \textbf{Rule} & \textbf{Validation behaviour} \\
\midrule
\code{ticker}   & string  & Non-empty symbol             & UI prevents empty ticker \\
\code{quantity} & numeric & Positive or negative (signed) & Numeric handling enforced in data-entry flow \\
\bottomrule
\end{tabular}
\end{center}

\subsection{Option Input}

\begin{center}
\begin{tabular}{L{2.8cm}L{1.8cm}L{3.5cm}L{4.5cm}}
\toprule
\textbf{Field} & \textbf{Type} & \textbf{Rule} & \textbf{Validation behaviour} \\
\midrule
\code{label}/\code{ticker}    & string & Non-empty & UI avoids blank partial rows \\
\code{underlying\_ticker}      & string & Must exist in price data & Explicit ticker-existence check implemented \\
\code{option\_type}           & string & \code{call} or \code{put} & Pricing layer rejects unknown types \\
\code{quantity}               & numeric & Signed & Numeric handling enforced \\
\code{strike}                 & numeric & Positive & Invalid values fail in pricing \\
\code{maturity}               & date    & Future-dated for live pricing & Expired options handled; malformed dates fail upstream \\
\code{volatility}             & numeric & Positive decimal & Pricing layer rejects non-positive volatility \\
\code{risk\_free\_rate}       & numeric & Numeric decimal & Numeric input expected \\
\code{dividend\_yield}        & numeric & Numeric decimal & Numeric input expected \\
\code{multiplier}             & numeric & Positive & Expected positive; no central validator for all paths \\
\bottomrule
\end{tabular}
\end{center}

\subsection{Risk Settings}

\begin{center}
\begin{tabular}{L{2.8cm}L{1.8cm}L{3cm}L{4.5cm}}
\toprule
\textbf{Field} & \textbf{Type} & \textbf{Rule} & \textbf{Validation behaviour} \\
\midrule
Lookback window   & integer & Positive, sufficiently large & Infeasible windows surface as insufficient-history error \\
Horizon           & integer & Positive                    & Return helpers reject invalid horizons \\
VaR confidence    & float   & $(0,1)$                     & Domain covered in tests \\
ES confidence     & float   & $(0,1)$                     & Separate-confidence tests exist \\
Estimator type    & enum    & \code{window} or \code{ewma} & UI constrains; orchestration assumes supported values \\
Calibration mode  & enum    & \code{historical} or \code{manual} & UI constrains; tests cover manual mode \\
MC simulations    & integer & Positive                    & Positive value expected \\
Random seed       & integer or \code{None} & Fixed for reproducibility & Both fixed and unfixed usage supported \\
Vol shock mode    & enum    & \code{fixed} or \code{underlying\_beta} & UI constrains; unknown modes raise visibly \\
\bottomrule
\end{tabular}
\end{center}

\subsection{Output Schema}

\begin{center}
\begin{tabular}{L{3.5cm}L{9cm}}
\toprule
\textbf{Output} & \textbf{Contents} \\
\midrule
Risk summary (JSON) & Method, VaR, ES, VaR/ES confidence levels, horizon, portfolio value, assumptions \\
Losses CSV          & Scenario id or date, method, loss under convention $V_0 - V_T$ \\
Backtest CSV        & Model, observations, exceptions, expected exceptions, exception rate, Kupiec LR, p-value \\
\bottomrule
\end{tabular}
\end{center}

% ─────────────────────────────────────────────────────────────────────────────
\section{Risk Engine Orchestration}

The core orchestration path is implemented in
\code{src/services/risk\_engine\_service.py}.

\subsection{Core Steps}

\begin{enumerate}
  \item Accept a validated \code{Portfolio}, aligned \code{prices} DataFrame,
    a \code{pricing\_date}, and a risk-settings bundle.
  \item Compute current portfolio value from the latest spot vector.
  \item Dispatch to \code{historical\_var\_es}, \code{parametric\_var\_es},
    and \code{monte\_carlo\_var\_es}.
  \item Aggregate results into a single dictionary keyed by method.
  \item On backtest requests, call \code{run\_backtest} and then
    \code{kupiec\_test}.
  \item Return service-level objects to the UI for display and download.
\end{enumerate}

The settings bundle carries: \code{calibration\_mode}, optional
\code{manual\_market\_params} for mean/covariance overrides,
\code{option\_vol\_shock\_mode}, \code{option\_vol\_shock\_beta}, and
\code{option\_vol\_shock\_floor}.

\subsection{Why the Service Layer Matters}

The service layer isolates each Streamlit panel from knowledge of
portfolio repricing, return estimation, scenario generation, output
aggregation, and error conditions.  Embedding those responsibilities
in the UI would increase duplication and coupling.  Acting as a single
orchestration boundary makes the application more maintainable and its
model governance more tractable.

\subsection{Extension Orchestration}

Credit and regulatory logic use separate service modules
(\code{src/services/credit\_service.py} and
\code{src/services/regulatory\_service.py}), which keeps the required
stock/option risk workflow independent from the course-formula extensions.

% ─────────────────────────────────────────────────────────────────────────────
\section{Data Validation and Error Handling}

The codebase uses both front-end and back-end validation.  UI-level
controls reduce obvious user errors, while data- and model-layer checks
protect against invalid numerical states.

\subsection{Error-Handling Table}

\begin{center}
\begin{tabular}{L{4.5cm}L{3.5cm}L{5.5cm}}
\toprule
\textbf{Error case} & \textbf{Detection layer} & \textbf{Required behaviour} \\
\midrule
Missing ticker in price history  & Data validation & Raise or display explicit error \\
Too few observations              & Risk module & Return explicit empty reason in backtest paths; fail visibly in calculation paths \\
Negative or zero prices           & Data validation & Reject explicitly \\
Duplicate dates                   & Data loader & Reject explicitly; not all cases are centrally normalized \\
NaN prices                        & Data loader/validation & All-NaN columns flagged; broader handling is workflow-dependent \\
Invalid option maturity           & Schema/pricing & Expired options handled; malformed values fail visibly \\
Negative volatility               & Schema/pricing & Pricing fails visibly when repricing attempted \\
Invalid confidence                & Risk settings & Constrained by UI or fails visibly if misused \\
Non-PSD covariance                & Estimator/MC & Manual inputs rejected explicitly; auto-repair not the primary path \\
Empty portfolio                   & Schema/service & Prevented or fails visibly \\
Download failure                  & Data layer/UI & Clear user-facing error message \\
Monte Carlo seed missing          & MC layer & Randomize and record, or require seed for regression \\
\bottomrule
\end{tabular}
\end{center}

\subsection{Data Validation Design}

Lecture~5 emphasizes documenting data quality, proxies, and cleaning
assumptions.  We address these requirements at three layers:

\begin{itemize}
  \item \code{src/data/validation.py} checks emptiness, index type,
    all-NaN columns, positive prices, and stale-price runs.
  \item \code{src/data/market\_data.py} handles CSV parsing, numeric
    coercion, Yahoo Finance retrieval, exponential-backoff retry, and
    parquet caching.
  \item \code{src/ui/market\_data\_panel.py} surfaces errors immediately
    to the user rather than silently proceeding.
\end{itemize}

Visible failure behavior is mandatory for a model-risk-sensitive
application.  Enforcement is distributed across the UI, data loaders,
validators, and numerical modules; this document distinguishes the
intended input contract from the specific layer where each check is applied.

% ─────────────────────────────────────────────────────────────────────────────
\section{Numerical Implementation Controls}

Floating-point outputs are compared with tolerances rather than exact
equality throughout the test suite.  Domains are validated before
formulas are evaluated.  Monte Carlo regression tests use fixed seeds.
Covariance matrices are checked for shape, symmetry, and positive
semi-definiteness via estimator tests, and edge cases are exercised to
confirm finite outputs or controlled failures.

\begin{center}
\begin{tabular}{L{6cm}L{7.5cm}}
\toprule
\textbf{Numerical risk} & \textbf{Control} \\
\midrule
Floating-point equality             & \code{pytest.approx} / \code{np.isclose} in all tests \\
Overflow in lognormal shocks        & Validate inputs; use analytical formulas where available \\
Underflow in deep OTM option prices & Permit small positive values; reject negative option prices \\
Catastrophic cancellation           & Prefer analytical Greeks over unstable finite differences \\
Non-PSD covariance                  & Test structure and document error or repair expectations \\
Monte Carlo randomness              & Fixed seed for all regression-style tests \\
Stale or missing data               & Validate before calculation \\
Wrong loss sign                     & Explicit PnL/loss conventions plus sign-aware tests \\
VaR/ES confidence confusion         & Separate-confidence tests in \code{test\_es\_confidence\_split.py} \\
Delta/exposure unit mismatch        & Dedicated regression test confirms delta-dollar convention with spot \\
\bottomrule
\end{tabular}
\end{center}

% ─────────────────────────────────────────────────────────────────────────────
\section{Testing and Coverage}

Testing is integral to the design of the repository, not a sidecar
to the application.

\subsection{Test Integration by Layer}

\begin{itemize}
  \item Formula modules are validated by analytical and course-derived tests.
  \item Service orchestration is validated by backend integration tests.
  \item UI panels are validated with Streamlit panel tests.
  \item Market-data wrappers are validated independently from the risk engine.
  \item Integration scripts exercise end-to-end behavior with live market data.
\end{itemize}

\subsection{Test Execution Summary}

No-network unit test command:
\begin{shellcode}
python -m pytest tests/ \
  --ignore=tests/integration_test.py \
  --ignore=tests/integration_test_formula_sheet.py -v
\end{shellcode}

Recorded result (2026-05-11 03:00:09 EDT, commit \code{5841589e}):
\textbf{624 passed, 0 failed, 0 skipped}.

Coverage command:
\begin{shellcode}
python -m pytest tests/ --cov=src --cov-report=term-missing \
  --cov-report=html:submission/coverage_report \
  --cov-report=xml:submission/coverage_report/coverage.xml \
  --ignore=tests/integration_test.py \
  --ignore=tests/integration_test_formula_sheet.py
\end{shellcode}

Recorded result: \textbf{95\% statement coverage}.  Remaining untested
lines are concentrated in UI branch paths, selected credit-service
helpers, and defensive validation branches.

\subsection{Coverage and Remediation Plan}

The model-critical paths (pricing formulas, VaR/ES formulas, covariance
handling, backtesting diagnostics, credit formulas, and regulatory
arithmetic) are all tested directly.

\begin{center}
\begin{tabular}{L{3.5cm}L{5.5cm}L{4.5cm}}
\toprule
\textbf{Area} & \textbf{Remaining gap} & \textbf{Planned remediation} \\
\midrule
Streamlit panel branches  & Some UI display paths require live Streamlit/browser state & Add headless browser or Streamlit component tests as a future enhancement \\
Credit-service helpers    & Some orchestration branches are lower value than pure formula branches & Add service-level synthetic workflows \\
Defensive validation      & Some invalid states are hard to trigger from the UI & Add direct unit tests against validators \\
Historical edge paths     & Some insufficient-history and absolute-shock branches are method-specific & Add synthetic price matrices with controlled missing or short history \\
\bottomrule
\end{tabular}
\end{center}

% ─────────────────────────────────────────────────────────────────────────────
\section{Deployment and Reproducibility}

\subsection{Core Commands}

Install dependencies:
\begin{shellcode}
pip install -r requirements.txt
\end{shellcode}

Run the application:
\begin{shellcode}
streamlit run app.py
\end{shellcode}

Run the no-network test suite:
\begin{shellcode}
python -m pytest tests/ \
  --ignore=tests/integration_test.py \
  --ignore=tests/integration_test_formula_sheet.py
\end{shellcode}

\subsection{Recorded Test Environment}

\begin{itemize}
  \item Date/time: 2026-05-11 03:00:09 EDT
  \item Git commit: \code{5841589e3f3d2dbd3c1e38b08642eccce201a6a2}
  \item Python: 3.12.2 \quad OS: Darwin 24.5.0 arm64
  \item Key packages: \code{streamlit 1.37.1}, \code{numpy 1.26.4},
    \code{pandas 3.0.2}, \code{scipy 1.17.1}, \code{plotly 5.24.1},
    \code{yfinance 1.2.0}, \code{pytest 7.4.4}, \code{pytest-cov 7.1.0}
\end{itemize}

\subsection{Reproducibility Assessment}

Reproducibility is strong for deterministic and no-network paths:
analytical modules are deterministic, Monte Carlo defaults to a fixed
seed for regression-stable tests, and the application depends on
explicit settings objects.  Reproducibility is inherently weaker for
live-data integration because Yahoo Finance data may update, market
data availability can change, and external dependencies may vary.

% ─────────────────────────────────────────────────────────────────────────────
\section{Known Software Limitations}

\begin{center}
\begin{tabular}{L{4cm}L{4.5cm}L{5cm}}
\toprule
\textbf{Limitation} & \textbf{Consequence} & \textbf{Mitigation} \\
\midrule
Streamlit app is not production software & No enterprise auth, audit, or access controls & Academic use only; limitations are explicitly documented \\
Yahoo Finance / yfinance data can be imperfect & Corporate actions, stale prices, or connection failures may affect live runs & Retry logic and parquet cache implemented; Bloomberg CSV data available as fallback \\
Option pricing is Black-Scholes only & No stochastic volatility, no smile modeling & Appropriate for the course scope; documented as a known limitation \\
Parametric VaR uses delta-normal linear approximation & Nonlinear option payoffs are underestimated relative to full repricing & Historical and MC engines use full repricing; limitation is disclosed in model documentation \\
EWMA parameter is formula-derived, not ML-calibrated & May not match market-observed volatility clustering & Acceptable for academic purpose; documented in model description \\
Regulatory extensions are illustrative & Do not claim official Basel or DFAST compliance & Academic demonstration only; labeled as course formula implementations \\
Monte Carlo convergence is seed-dependent & Seeded runs are deterministic; unseeded runs introduce variance & Regression tests use fixed seeds; unseeded usage is documented \\
95\% statement coverage rather than 100\% & Some UI branches and defensive paths are untested & Remaining gaps are non-formula paths; documented in Section~10.3 \\
\bottomrule
\end{tabular}
\end{center}

% ─────────────────────────────────────────────────────────────────────────────
\section{Design Compliance with Lecture 5}

Lecture~5 frames pre-deployment model risk management in terms of five
requirements: requirements definition, design documentation, data
analysis, testing, and system analysis and testing.  We address each
explicitly:

\begin{enumerate}
  \item \textbf{Requirements.}  The system purpose and scope are stated
    in Section~2, including in-scope and out-of-scope boundaries.
  \item \textbf{Design documentation.}  Sections 3–5 provide a layered
    architecture diagram, separation-of-concerns table, module inventory,
    public interface contracts, and core data structures.
  \item \textbf{Data analysis.}  Section~8 documents the data validation
    strategy; Section~6 documents the input schema and data quality checks.
  \item \textbf{Testing.}  Section~10 documents the no-network test suite
    (624 tests, 95\% coverage), integration scripts, and coverage
    remediation plan.
  \item \textbf{System analysis and testing.}  Sections 4 and 7 describe
    the end-to-end data flow and the service-layer orchestration that binds
    the system together.
\end{enumerate}

Lecture~5 also explicitly calls out bugs, design weaknesses, hidden
assumptions, floating-point comparison errors, unit tests, single-purpose
routines, and separation of data gathering from calibration from
computation.  Sections 8–12 address each of these concerns directly.

\end{document}
